% 2009 年天津卷文科第 19 题：四棱锥 P-ABCD，PD⊥平面 ABCD，E 为 PC 中点。
% 来源 PDF：普通高考/2009/2009天津文.pdf，SHA256 5d3e0110cc9145311eed3af1cf0f691ee327cbcba3a0d245e88d635fc136203f
% 原页：第 2 页（A3 三栏）第 1 栏，印刷页 564。
% 重绘依据：按原图 P 在 D 上方，D 左下、C 右下、B 最右、A 在 D 上方偏左；
%   棱 PA、PB、PC 实线，底面边 AB、BC、CD、DA 与对角线 DB 实线，AC 画虚线；
%   E 为 PC 中点并标出。
\begin{tikzpicture}[scale=1.0, semithick, every node/.style={inner sep=1.2pt}]
  \coordinate (P) at (0.5,3.6);
  \coordinate (D) at (0,0);
  \coordinate (C) at (2.4,-0.3);
  \coordinate (B) at (4.2,1.3);
  \coordinate (A) at (0.9,1.6);
  \coordinate (E) at ($(P)!0.5!(C)$);
  % 底面
  \draw (A) -- (B) -- (C) -- (D) -- cycle;
  \draw (D) -- (B);
  \draw[dashed] (A) -- (C);
  % 侧棱
  \draw (P) -- (A);
  \draw (P) -- (B);
  \draw (P) -- (C);
  % E
  \fill (E) circle (1.2pt);
  \node[anchor=south] at (P) {$P$};
  \node[anchor=north east] at (D) {$D$};
  \node[anchor=north west] at (C) {$C$};
  \node[anchor=west] at (B) {$B$};
  \node[anchor=east] at (A) {$A$};
  \node[anchor=west] at ([xshift=3pt]E) {$E$};
\end{tikzpicture}
